% Copie este arquivo para uma ficha numerada e preencha somente informações verificadas.
% Inclua a nova ficha em main.tex, separada das demais por \clearpage.
\fichatitulo{00 · EIXO TEMÁTICO}{Tipo de publicação · Ano}{Título curto}
{\small Referência completa e link para a fonte consultada.\par}

\campo{Problema e objetivo}
Questão enfrentada pelo texto, com suas próprias palavras.

\campo{Principal contribuição · síntese interpretativa}
Contribuição específica, distinguindo proposta e resultado demonstrado.

\campo{Método e evidência}
Desenho, dados e evidências. Para textos institucionais, explicitar estrutura e recorte.

\campo{Resumo extrativo · fragmentos literais selecionados}
\begin{itemize}
\excerto{Ideia central}{Trecho literal curto conferido na fonte}{Página e seção da versão consultada.}
\excerto{Segunda ideia}{Trecho literal curto conferido na fonte}{Página e seção da versão consultada.}
\excerto{Terceira ideia}{Trecho literal curto conferido na fonte}{Página e seção da versão consultada.}
\end{itemize}

\campo{Limitações e alcance}
Separar limitações declaradas pelo texto e observações críticas do fichamento.

\campo{Relação com a dissertação · análise proposta}
Indicar onde e como a referência pode ser usada, sem antecipar decisões do pesquisador.

\registro{Versão e escopo efetivamente consultados nesta preparação.}
